\chapter{Queue-Reactive Models and Large-Tick Assets}\label{ch:hf:queue-reactive-models-and-large-tick-assets}

On a stock whose spread is one tick nearly all day, every market maker quotes the same two prices. What separates them is where they stand in the
queue, and when they leave it. On this chapter's simulated stock the spread is one tick 96\% of the time. A one-lot quoter that always joins the
best queues earns 0.03 ticks a share over the ten seconds after its fills; one that joins only the side whose queue is at least 30\% of the two,
and leaves when that share falls, earns 0.39 on a third of the volume. The same join rule without the exit loses 0.12. Leaving is the part that
pays.

\section{The large-tick regime}

One Quant Book 10, chapter 6 classifies instruments by their relative tick size: in a \emph{large-tick asset} the spread is almost always one tick
and the best queues are long, in a \emph{small-tick asset} the spread spans many ticks and queues are short. The distinction decides what a market
maker controls. In a small-tick asset the quote's distance from the mid is a continuous choice, the depth of chapter 3. In a large-tick asset the
price is fixed by the tick grid and the only choices are to be in a best queue or not, and where in it.

Book 7's simulated stock is large-tick. Over an hour, its spread is one tick 96.3\% of the time and the median best queue holds 17 lots of 100
shares. Chapter 3's exploring quoter found no fill at all a tick behind the best price. Everything a market maker earns there, it earns in the two
best queues.

\section{Queue-reactive intensities}

\begin{definition}[Queue-reactive model]\label{def:hf:queue-reactive-models-and-large-tick-assets:qr}
A \emph{queue-reactive model}\index{queue-reactive model} describes each best queue of a book as a birth--death process whose rates depend on the
queue's current size $q$: limit orders join its back at rate $\lambda^L(q)$, cancellations remove resting orders at rate $\lambda^C(q)$, market
orders remove orders from its front at rate $\lambda^M(q)$. When a queue empties the price moves by one tick and new queues are drawn from a
regeneration distribution.
\end{definition}

Huang, Lehalle and Rosenbaum introduced the model on real order-book data: between changes of a reference price the book is a Markov queueing
system whose intensities depend only on its current state. It is the birth--death chain of One Quant Book 4, chapter 8, with rates read off the
data instead of assumed constant, as in Cont, Stoikov and Talreja's earlier model.

\begin{method}[Estimating the intensities]\label{met:hf:queue-reactive-models-and-large-tick-assets:estimate}
Replay the order-by-order messages. For each message, note the size of each best queue just before it and the time since the previous message;
credit that time to the size. Count the lots added to, cancelled from and executed against each best queue at each size. The intensity at size $q$
is the count divided by the time spent at $q$; pool the two sides if the book is symmetric.
\end{method}

\Cref{fig:hf:queue-reactive-models-and-large-tick-assets:intensities} shows the estimate on one simulated hour. Cancellations grow with the queue,
from 0.34 lots a second at one lot to about 1.5 at ten, because each resting order cancels at its own rate; market orders reach small queues more
often than large ones (0.55 a second at two lots, about 0.3 at fifteen). Limit orders arrive at 1.3 to 2.3 lots a second, with no trend in the size.

\begin{omfigure}
\begin{tikzpicture}
\begin{axis}[width=11cm,height=5.2cm,xlabel={\small queue size before the event (lots)},ylabel={\small lots per second},
  tick label style={font=\footnotesize},xmin=1,xmax=25,ymin=0,ymax=2.6,grid=major,grid style={gray!25},table/col sep=comma,
  legend style={font=\scriptsize,at={(0.5,-0.3)},anchor=north},legend columns=3]
\addplot[omDef,thick,mark=*,mark size=1pt] table[x=q,y=L]{figdata/hft/05-queue-reactive-models-and-large-tick-assets/intensities.csv};
\addplot[omThm,thick,dashed,mark=square*,mark size=1pt] table[x=q,y=C]{figdata/hft/05-queue-reactive-models-and-large-tick-assets/intensities.csv};
\addplot[black,thick,mark=triangle*,mark size=1.2pt] table[x=q,y=M]{figdata/hft/05-queue-reactive-models-and-large-tick-assets/intensities.csv};
\legend{limit orders $\lambda^L$,cancellations $\lambda^C$,market orders $\lambda^M$}
\end{axis}
\end{tikzpicture}

\omcaption{Queue-reactive intensities of the best queues of the simulated stock, pooled over bid and ask, by the queue's size just before each
event; one simulated hour. Data: \texttt{firm.qreactive.estimate} on \texttt{firm.tape}.}\label{fig:hf:queue-reactive-models-and-large-tick-assets:intensities}
\end{omfigure}

\section{The value of a place in the queue}

One Quant Book 10, chapter 6 defines the \emph{queue value} of a resting order: what the order is expected to earn from its current place until
it is filled or cancelled. Moallemi and Yuan split it into a static part, the spread earned against adverse selection costs that grow with the
position in the queue, and a dynamic part, the option to leave later; for some large-tick stocks they found it of the same order as the spread.

The fitted \omterm{def:hf:queue-reactive-models-and-large-tick-assets:qr}{queue-reactive model} computes it by simulation. Put one lot of ours in the bid queue with $n$ lots ahead; the bid queue holds $Q^b$
lots in all and the ask queue $Q^a$. Market sells take the bid from the front, cancellations hit the other orders at random, limit orders join the
back. The order is filled when a market sell reaches it; it is abandoned, worth zero, if the ask queue empties first (the price moves up, away
from it) or after sixty seconds. After a fill the mid is followed for ten seconds, each emptied queue moving it by a tick, and the order is marked
there: half a tick plus the mid's move (\cref{lst:hf:queue-reactive-models-and-large-tick-assets:mc}).

\begin{omfigure}
\begin{tikzpicture}
\begin{axis}[width=11cm,height=5.2cm,xlabel={\small size of our side's queue $Q^b$ (lots)},ylabel={\small value (ticks a share)},xmode=log,
  log basis x=10,xtick={2,5,10,20},xticklabels={2,5,10,20},xmin=1.8,xmax=22,ymin=0,ymax=0.7,tick label style={font=\footnotesize},
  grid=major,grid style={gray!25},table/col sep=comma,legend style={font=\scriptsize,at={(0.5,-0.3)},anchor=north},legend columns=2]
\addplot[omDef,thick,mark=*,mark size=1.2pt] table[x=same,y=front2]{figdata/hft/05-queue-reactive-models-and-large-tick-assets/values.csv};
\addplot[omDef,thick,dashed,mark=o,mark size=1.2pt] table[x=same,y=back2]{figdata/hft/05-queue-reactive-models-and-large-tick-assets/values.csv};
\addplot[omThm,thick,mark=square*,mark size=1.2pt] table[x=same,y=front20]{figdata/hft/05-queue-reactive-models-and-large-tick-assets/values.csv};
\addplot[omThm,thick,dashed,mark=square,mark size=1.2pt] table[x=same,y=back20]{figdata/hft/05-queue-reactive-models-and-large-tick-assets/values.csv};
\legend{{front, $Q^a=2$},{back, $Q^a=2$},{front, $Q^a=20$},{back, $Q^a=20$}}
\end{axis}
\end{tikzpicture}

\omcaption{Value of one lot resting in the bid queue, at its front and at its back, in the fitted \omterm{def:hf:queue-reactive-models-and-large-tick-assets:qr}{queue-reactive model}, by the bid queue's size,
for a short and a long ask queue: expected P\&L per share marked ten seconds after the fill (zero if not filled within sixty seconds or if the
ask queue empties first); 4\,000 simulated paths per point. Data: \texttt{hf\_queues.values}.}\label{fig:hf:queue-reactive-models-and-large-tick-assets:values}
\end{omfigure}

Three regularities show in \cref{fig:hf:queue-reactive-models-and-large-tick-assets:values}. The front is worth more than the back when our
queue is long: at ten lots with a short ask queue, 0.60 against 0.32 ticks. When our queue is short and the other long, front and back are worth
little and almost the same (0.12 at two lots against twenty): a fill is nearly certain because the queue is about to be eaten, and the price follows.
When the other queue is short, a back order risks never being filled, because the price moves away first.

\begin{remark}[What the model cannot see]\label{rem:hf:queue-reactive-models-and-large-tick-assets:info}
Every value in the figure is positive, yet chapter 1's quoter, always in these queues, lost its spread to adverse selection. The fitted model
knows only queue sizes. In the simulated market, as in real ones, some market orders come from traders who know where the efficient price is,
and they arrive when the quotes are stale. A \omterm{def:hf:queue-reactive-models-and-large-tick-assets:qr}{queue-reactive model} captures the mechanical race for the front; the information in the flow must be
measured on one's own fills, by mark-outs, and added to it (chapter 6).
\end{remark}

\section{When to join and when to leave}

\begin{definition}[Queue-join rule, queue-exit rule]\label{def:hf:queue-reactive-models-and-large-tick-assets:rules}
A \emph{queue-join rule}\index{queue-join rule} decides whether to place an order at the back of a best queue, given the state of the book; a
\emph{queue-exit rule}\index{queue-exit rule} decides whether to cancel an order already resting, given the state of the book and the order's place
in its queue.
\end{definition}

The two rules are different decisions: joining forfeits nothing, leaving forfeits the place. The chapter's rules use one number, the share of our
side's best queue in the two best queues, counting others' orders only, $s^b=Q^b/(Q^b+Q^a)$ for the bid (and $1-s^b$ for the ask). The model says a
side whose share is small is about to be eaten: fills there are nearly certain and followed by an adverse move. A one-lot quoter joins a side
only when its share is at least $\theta$; the exit rule then either stays, leaves when the share falls below $\theta$, or leaves unless fewer than
two lots are ahead (\cref{lst:hf:queue-reactive-models-and-large-tick-assets:quoter}). On six simulated sessions of twenty minutes:

\begin{center}
\small
\begin{tabular}{@{}lrrrr@{}}
\toprule
rule ($\theta=0.3$) & shares & ten-second mark-out & ten-second edge & P\&L (sd)\\
& a session & (ticks a share) & (\$ a session) & (\$ a session)\\
\midrule
always join & 11\,167 & 0.026 & 2.92 & 5.33 (63.67)\\
join, then stay & 6\,900 & $-0.117$ & $-8.08$ & $-32.00$ (68.41)\\
join and leave & 3\,400 & 0.392 & 13.33 & $-6.50$ (23.82)\\
join, leave unless at the front & 4\,450 & 0.032 & 1.42 & $-26.50$ (63.22)\\
\bottomrule
\end{tabular}
\end{center}

The ten-second mark-out, spread capture less the adverse move over ten seconds, is the fill's edge; times the shares, it is the session's edge. By
that measure the exit rule is what works. Joining on a favourable share and then staying is worse than always joining: the orders that stay are
exactly those whose side turned against them, and they are the ones that get filled. Leaving when the share falls quadruples the edge on a third of
the volume. Keeping orders near the front, the natural refinement, gives most of the gain back, as the model predicted: when the other side is
about to take the queue, being first means being filled first. The P\&L column, which also marks the inventory to the end of each session, is too
noisy over six sessions to rank the rules.

\Cref{fig:hf:queue-reactive-models-and-large-tick-assets:sweep} varies $\theta$ for the join-and-leave rule. The mark-out per share rises with the
threshold up to 0.4; the edge per session peaks at 0.3, where the quoter still trades enough.

\begin{omfigure}
\begin{tikzpicture}
\begin{axis}[width=11cm,height=5.2cm,axis y line*=right,axis x line=none,ylabel={\small ten-second mark-out (ticks a share)},
  tick label style={font=\footnotesize},xmin=-0.02,xmax=0.52,ymin=0,ymax=0.6,table/col sep=comma]
\addplot[omDef,thick,dashed,mark=square*,mark size=1.3pt] table[x=theta,y=markout10]{figdata/hft/05-queue-reactive-models-and-large-tick-assets/sweep.csv};
\label{plot:hf:qsweep}
\end{axis}
\begin{axis}[width=11cm,height=5.2cm,axis y line*=left,xlabel={\small threshold $\theta$ on our side's share of the best queues},
  ylabel={\small ten-second edge (\$ a session)},tick label style={font=\footnotesize},xmin=-0.02,xmax=0.52,ymin=0,ymax=15,grid=major,
  grid style={gray!25},table/col sep=comma,legend style={font=\scriptsize,at={(0.5,-0.3)},anchor=north},legend columns=2]
\addplot[omThm,thick,mark=*,mark size=1.3pt] table[x=theta,y=edge10]{figdata/hft/05-queue-reactive-models-and-large-tick-assets/sweep.csv};
\addlegendentry{edge a session (left)}
\addlegendimage{/pgfplots/refstyle=plot:hf:qsweep}\addlegendentry{mark-out a share (right)}
\end{axis}
\end{tikzpicture}

\omcaption{The join-and-leave rule on six simulated sessions of twenty minutes: ten-second edge a session and ten-second mark-out per share filled,
against the threshold $\theta$ ($\theta=0$ is always joining). Data: \texttt{hf\_queues.sweep}.}\label{fig:hf:queue-reactive-models-and-large-tick-assets:sweep}
\end{omfigure}

\section{Tutorial: pricing a place in the queue}\label{tut:hf:queue-reactive-models-and-large-tick-assets:tut}

\begin{tutorial}
\textbf{Goal.} Estimate the queue-reactive intensities of a large-tick book, value a place in its queue, and test join and exit rules in the
market. \textbf{End state:} \cref{fig:hf:queue-reactive-models-and-large-tick-assets:intensities,fig:hf:queue-reactive-models-and-large-tick-assets:values,fig:hf:queue-reactive-models-and-large-tick-assets:sweep}
and the table.
\begin{tutsteps}
\item \textbf{Estimate} with \texttt{firm.qreactive.estimate(msgs, top, lot, qmax)} on one simulated hour (\cref{met:hf:queue-reactive-models-and-large-tick-assets:estimate}).
\item \textbf{Value the queue} with \texttt{order\_value}: all states and paths advance together, one event each per step. The bid side's events
are the heart of it.
\omcode{code/firm/qreactive/firm_qreactive.py}{121}{138}{Our order in the bid queue: joins behind it, cancellations ahead of it, market orders from the front.}\label{lst:hf:queue-reactive-models-and-large-tick-assets:mc}
\item \textbf{Rules as a Quoter.} The harness gives the quoter its order's place in the queue (\texttt{ctx.ahead}), as an order-by-order feed
would; the quoter reads the others' queues from \texttt{ctx.external}.
\omcode{code/hft/05-queue-reactive-models-and-large-tick-assets/python/hf_queues.py}{79}{92}{Join a side on its share of the best queues; stay, leave, or leave unless at the front.}\label{lst:hf:queue-reactive-models-and-large-tick-assets:quoter}
\item \textbf{Compare} with \texttt{compare()} and \texttt{sweep()} on sessions never used for estimation.
\end{tutsteps}
\textbf{What to change next.} Add the imbalance of the next level; let the threshold depend on the place in the queue, using the model's value.
\end{tutorial}

\section{Build: the queue-reactive engine}\label{bld:hf:queue-reactive-models-and-large-tick-assets:qreactive}

\begin{build}
\bfield{Purpose} Queue-reactive intensities from any order-by-order stream, and the value of an order's place in its queue, for join and exit
decisions.
\bfield{Interface} \texttt{estimate(msgs, top, lot, qmax, n\_open) -> \{q, L, C, M, time, regen\}}; \texttt{QRModel(L, C, M, regen)};
\texttt{order\_value(model, ahead, same, opp, H, tmax, paths, seed) -> \{fill, value, adverse\}}; \texttt{value\_table}. In
\texttt{firm.mmharness}, \texttt{ctx.ahead(cid)} for Quoters that set \texttt{wants\_queue\_position}.
\bfield{Rules} Events are counted in lots, one lot per event. Estimation and evaluation never use the same sessions. Values are per share, in
ticks, marked at the mid a fixed horizon after the fill.
\bfield{Acceptance tests} \texttt{code/firm/qreactive/tests/}: intensities and time in state by hand on a four-message stream; no fill without market
orders; the front fills more often than the back; table shapes. \texttt{code/firm/mmharness/tests/}: an order's place in its queue never moves back.
\bfield{Stretch} Intensities that depend on both queues (the full model); an information term fitted to mark-outs; run on \texttt{firm.exchsim}'s
feed.
\end{build}

\begin{omsources}
\item W.~Huang, C.-A.~Lehalle and M.~Rosenbaum, ``Simulating and analyzing order book data: the queue-reactive model'', \emph{Journal of the American
Statistical Association} 110(509), 2015.
\item C.~C.~Moallemi and K.~Yuan, ``A model for queue position valuation in a limit order book'', working paper, 2016.
\item R.~Cont, S.~Stoikov and R.~Talreja, ``A stochastic model for order book dynamics'', \emph{Operations Research} 58(3), 2010.
\item K.~Dayri and M.~Rosenbaum, ``Large tick assets: implicit spread and optimal tick size'', \emph{Market Microstructure and Liquidity} 1(1), 2015.
\end{omsources}

\section{Exercises}

\begin{exercise}[$\star$]\label{exo:hf:queue-reactive-models-and-large-tick-assets:1}
A best queue spent 200 seconds at 4 lots; in that time 290 lots were cancelled from it. Estimate $\lambda^C(4)$.
\end{exercise}

\begin{exercise}[$\star$]\label{exo:hf:queue-reactive-models-and-large-tick-assets:2}
The bid queue holds 3 lots of others' orders and the ask queue 12. What is the bid's share, and does the chapter's rule with $\theta=0.3$ join the
bid? The ask?
\end{exercise}

\begin{exercise}[$\star$]\label{exo:hf:queue-reactive-models-and-large-tick-assets:3}
A quoter's fills earned 0.392 ticks a share after ten seconds on 3\,400 shares a session. With a one-cent tick, what is its ten-second edge a
session?
\end{exercise}

\begin{exercise}[$\star\star$]\label{exo:hf:queue-reactive-models-and-large-tick-assets:4}
Why can a place at the back of a queue be worth more when the opposite queue is long than when it is short?
\end{exercise}

\begin{exercise}[$\star\star$]\label{exo:hf:queue-reactive-models-and-large-tick-assets:5}
Explain why joining on a favourable share and then staying did worse than always joining.
\end{exercise}

\begin{exercise}[$\star\star$]\label{exo:hf:queue-reactive-models-and-large-tick-assets:6}
From \cref{fig:hf:queue-reactive-models-and-large-tick-assets:sweep}, why does the edge per session peak at a lower threshold than the mark-out
per share?
\end{exercise}

\begin{exercise}[$\star\star\star$]\label{exo:hf:queue-reactive-models-and-large-tick-assets:7}
\emph{Coding.} With the fitted model, compute the value, fill probability and adverse move of one lot at the front of a bid queue of 10 lots when
the ask queue holds 2 (\texttt{order\_value(model(), 0, 10, 2, paths=4000, seed=6)}).
\end{exercise}

\begin{exercise}[$\star\star\star$]\label{exo:hf:queue-reactive-models-and-large-tick-assets:8}
\emph{Find the flaw.} ``The \omterm{def:hf:queue-reactive-models-and-large-tick-assets:qr}{queue-reactive model} says every place in our queues is worth at least a tenth of a tick, so we should always quote.''
\end{exercise}

\section{Problem: Front of the Line}

\begin{problem}[{Weekend problem --- front of the line}]\label{pb:hf:queue-reactive-models-and-large-tick-assets:1}
A market maker in a one-tick stock must decide, message by message, which queues to be in.

\medskip\noindent\textbf{Part I --- The regime.}
\begin{enumerate}
\item What distinguishes a large-tick asset, and what can a market maker control there?
\item Give the simulated stock's one-tick share and median best queue.
\item Why did chapter 3's exploring quoter get no fill behind the touch?
\end{enumerate}

\medskip\noindent\textbf{Part II --- The model.}
\begin{enumerate}[resume]
\item Define the \omterm{def:hf:queue-reactive-models-and-large-tick-assets:qr}{queue-reactive model}.
\item Describe the estimation and give the fitted intensities' shapes.
\item Why do cancellations grow with the queue in this market?
\item Describe the simulation that values a place in the queue, and its payoff.
\end{enumerate}

\medskip\noindent\textbf{Part III --- The value of a place.}
\begin{enumerate}[resume]
\item Give the front and back values at 10 lots against 2, and at 2 lots against 20.
\item Why is the front not worth more than the back when our queue is short and the other long?
\item Why are all the model's values positive, and what does that omit?
\end{enumerate}

\medskip\noindent\textbf{Part IV --- Rules in the market.}
\begin{enumerate}[resume]
\item Define the \omterm{def:hf:queue-reactive-models-and-large-tick-assets:rules}{queue-join rule} and the \omterm{def:hf:queue-reactive-models-and-large-tick-assets:rules}{queue-exit rule}.
\item State the \emph{named result}: the ten-second mark-out per share of always joining, joining and staying, and joining and leaving, and the value
of the front against the back of a ten-lot queue in the model.
\item Why is the P\&L column not used to rank the rules?
\item Why did keeping orders at the front give the gain back?
\item Which threshold maximises the edge per session, and why not the mark-out?
\item How many messages does each rule send, compared with always joining, and why?
\item How would you add an information term to the model?
\item What would change on a venue that allocates pro rata?
\item How would the rules behave on a small-tick asset?
\item In one sentence: what matters more in a one-tick book, joining well or leaving well?
\end{enumerate}
\end{problem}

\section{Interview questions}

\begin{interviewq}[$\star$ \iqroles{trader}]\label{iq:hf:queue-reactive-models-and-large-tick-assets:1}
You are first in a long bid queue in a one-tick stock. The ask queue has just shrunk to one lot. What do you do?
\end{interviewq}

\begin{interviewq}[$\star\star$ \iqroles{researcher}]\label{iq:hf:queue-reactive-models-and-large-tick-assets:2}
How would you estimate how much a place at the front of the queue is worth, from your own data?
\end{interviewq}

\begin{interviewq}[$\star\star$ \iqroles{researcher}]\label{iq:hf:queue-reactive-models-and-large-tick-assets:3}
Why do fill probability and adverse selection move together along a queue?
\end{interviewq}

\begin{interviewq}[$\star\star$ \iqroles{developer}]\label{iq:hf:queue-reactive-models-and-large-tick-assets:4}
How does a market maker know how many shares are ahead of its order, and what breaks the estimate?
\end{interviewq}

\begin{interviewq}[$\star\star$ \iqroles{trader}]\label{iq:hf:queue-reactive-models-and-large-tick-assets:5}
A colleague proposes to cancel and re-enter your orders every second to follow the book. What does that cost in a one-tick stock?
\end{interviewq}

\begin{interviewq}[$\star\star\star$ \iqroles{researcher}]\label{iq:hf:queue-reactive-models-and-large-tick-assets:6}
Each resting order cancels at rate $\nu$ and market orders arrive at rate $\mu$. An order has $n$ orders ahead of it. What is the probability that
the next event among those ahead is a market order, and what is the expected time until all $n$ are gone?
\end{interviewq}
